\documentclass[../DoAn.tex]{subfiles}
\begin{document}
Chương này có độ dài từ 3 đến  trang.

Với phương pháp phân tích thiết kế hướng đối tượng, sinh viên sử dụng biểu đồ use case theo hướng dẫn của template này. Với các phương pháp khác, sinh viên trao đổi với giáo viên hướng dẫn để đổi tên và sắp xếp lại đề mục cho phù hợp. Ví dụ, thay vì sử dụng biểu đồ use case, sinh viên đi theo hướng tiếp cận Agile có thể dùng User Story.

\section{Khảo sát hiện trạng}
\label{section:2.1}
Thông thường, khảo sát chi tiết về hiện trạng và yêu cầu của phần mềm sẽ được lấy từ ba nguồn chính, đó là (i) người dùng/khách hàng, (ii) các hệ thống đã có, (iii) và các ứng dụng tương tự.
Sinh viên cần tiến hành phân tích, so sánh, đánh giá chi tiết ưu nhược điểm của các sản phẩm/nghiên cứu hiện có. Sinh viên có thể lập bảng so sánh nếu cần thiết. Kết hợp với khảo sát người dùng/khách hàng (nếu có), sinh viên nêu và mô tả sơ lược các tính năng phần mềm quan trọng cần phát triển.

Đối với đề tài game...
This is a key subsection of your document. Here you must compare your game to others already developed. It is important to give a small description of the game being compared to, and point out the similarities between both. This is an excellent opportunity to expand the comparisons that were already made across the GDD and give the reader a better picture of what the game will actually be.

At the end, summarise your product's strong points and convince the reader why your product would sell despite its competitors. This is the trickiest part, because you must pick good opponents, otherwise the reader just won't know what you are talking about, and still keep your game’s image shining; therefore good writing is crucial. Your ‘adversaries’ also help with the notion of how big your market can be.

\section{Mục đích của trò chơi}
\label{section:2.2}
Pedagogical objective(s) for educational games, training goal, or intended social impact. If it is addressing a problem, explain what that problem is (for example, students have a hard time …)

\section{Định hướng sử dụng}
\label{section:2.3}
How the game is to be inserted into a pedagogical scenario or how will it be used in training, therapy, rehabilitation, etc. You should distinguish between a game that is stand alone, a supplement to a class, or working with an instructor

\section{Cơ sở của định hướng sử dụng}
\label{section:2.4}
Is there research that shows that this is apt to be successful? Are there are other games or applications that take a similar approach? If this is an experiment to see if it will work, why does the researcher believe that it might?

%%%%%%%%%%%%%%%%%%%%%%%%%%%%%%%%%%%

\end{document}